% theme: REFlesson_01hyou
\MajorQuestionLesson{連立方程式の意味と解：おかわり}
\SetRowSpace{5mm}

\TypeQuestionDouble{表の各列について、$x$ と $y$ の値の組が二元一次方程式を成り立たせるように、空欄を埋めなさい。}{hyou}

\begingroup
\newcommand{\SimTableCell}[1]{%
  \makebox[5mm][c]{\raisebox{0.3mm}{$#1$}}%
}
\renewcommand{\arraystretch}{1.4}

\QQRow{%
  $4x+3y=7$\par\smallskip
  \begin{tabular}{|c|c|c|c|c|c|}
    \hline
    \SimTableCell{x} & \SimTableCell{-2} & \SimTableCell{7} & \SimTableCell{}   & \SimTableCell{-5} & \SimTableCell{}    \\ \hline
    \SimTableCell{y} & \SimTableCell{}   & \SimTableCell{}  & \SimTableCell{13} & \SimTableCell{}   & \SimTableCell{-11} \\ \hline
  \end{tabular}%
}{左から $5,\ -7,\ -8,\ 9,\ 10$}{%
  $5x-2y=3$\par\smallskip
  \begin{tabular}{|c|c|c|c|c|c|}
    \hline
    \SimTableCell{x} & \SimTableCell{}    & \SimTableCell{3} & \SimTableCell{}   & \SimTableCell{}   & \SimTableCell{-1} \\ \hline
    \SimTableCell{y} & \SimTableCell{-14} & \SimTableCell{}  & \SimTableCell{11} & \SimTableCell{-9} & \SimTableCell{}   \\ \hline
  \end{tabular}%
}{左から $-5,\ 6,\ 5,\ -3,\ -4$}

\QQRow{%
  $2x+5y=13$\par\smallskip
  \begin{tabular}{|c|c|c|c|c|c|}
    \hline
    \SimTableCell{x} & \SimTableCell{9} & \SimTableCell{}  & \SimTableCell{-11} & \SimTableCell{}   & \SimTableCell{}  \\ \hline
    \SimTableCell{y} & \SimTableCell{}  & \SimTableCell{5} & \SimTableCell{}    & \SimTableCell{-3} & \SimTableCell{1} \\ \hline
  \end{tabular}%
}{左から $-1,\ -6,\ 7,\ 14,\ 4$}{%
  $3x-4y=-7$\par\smallskip
  \begin{tabular}{|c|c|c|c|c|c|}
    \hline
    \SimTableCell{x} & \SimTableCell{}   & \SimTableCell{}   & \SimTableCell{7} & \SimTableCell{-9} & \SimTableCell{}  \\ \hline
    \SimTableCell{y} & \SimTableCell{-8} & \SimTableCell{10} & \SimTableCell{}  & \SimTableCell{}   & \SimTableCell{4} \\ \hline
  \end{tabular}%
}{左から $-13,\ 11,\ 7,\ -5,\ 3$}

\QQRow{%
  $5x+y=-9$\par\smallskip
  \begin{tabular}{|c|c|c|c|c|c|}
    \hline
    \SimTableCell{x} & \SimTableCell{-4} & \SimTableCell{}  & \SimTableCell{}    & \SimTableCell{-1} & \SimTableCell{0}  \\ \hline
    \SimTableCell{y} & \SimTableCell{}   & \SimTableCell{1} & \SimTableCell{-14} & \SimTableCell{}   & \SimTableCell{}   \\ \hline
  \end{tabular}%
}{左から $11,\ -2,\ 1,\ -4,\ -9$}{%
  $x-5y=13$\par\smallskip
  \begin{tabular}{|c|c|c|c|c|c|}
    \hline
    \SimTableCell{x} & \SimTableCell{}   & \SimTableCell{-12} & \SimTableCell{-7} & \SimTableCell{}  & \SimTableCell{-2} \\ \hline
    \SimTableCell{y} & \SimTableCell{-1} & \SimTableCell{}    & \SimTableCell{}   & \SimTableCell{0} & \SimTableCell{}   \\ \hline
  \end{tabular}%
}{左から $8,\ -5,\ -4,\ 13,\ -3$}

\QQRow{%
  $4x-3y=-11$\par\smallskip
  \begin{tabular}{|c|c|c|c|c|c|}
    \hline
    \SimTableCell{x} & \SimTableCell{-14} & \SimTableCell{4} & \SimTableCell{}   & \SimTableCell{}  & \SimTableCell{1} \\ \hline
    \SimTableCell{y} & \SimTableCell{}    & \SimTableCell{}  & \SimTableCell{-7} & \SimTableCell{13} & \SimTableCell{}  \\ \hline
  \end{tabular}%
}{左から $-15,\ 9,\ -8,\ 7,\ 5$}{%
  $3x+5y=9$\par\smallskip
  \begin{tabular}{|c|c|c|c|c|c|}
    \hline
    \SimTableCell{x} & \SimTableCell{}  & \SimTableCell{8} & \SimTableCell{}   & \SimTableCell{-7} & \SimTableCell{}  \\ \hline
    \SimTableCell{y} & \SimTableCell{9} & \SimTableCell{}  & \SimTableCell{-6} & \SimTableCell{}   & \SimTableCell{0} \\ \hline
  \end{tabular}%
}{左から $-12,\ -3,\ 13,\ 6,\ 3$}

\endgroup
